%================================
% Vocabulary & Language Focus — XB1 (M–R)
%================================
\documentclass[../main.tex]{subfiles}

\begin{document}

%================================
\section{M}
%================================

\vw{metaphor}{/ˈmetəfər/}{n.}

\sense{n.}{暗喻，隐喻}
\col{a metaphor for sth（……的隐喻）}
\ex{The tunnel at the end of light is a common metaphor for despair.}
\ex{In the poem, the river is a metaphor for the passing of time.}
\col{a mixed metaphor（混搭失当的隐喻，修辞病）}
\ex{``We'll burn that bridge when we come to it'' mixes two metaphors.}

\begin{gram}{比喻三件套：metaphor / simile / personification}
\textbf{simile}（明喻）：用 \textit{like / as} 明言相似（\textit{brave as a lion}）；\textbf{metaphor}（暗喻）：不用比较词，直接以彼名此（\textit{Time is a river.}）；\textbf{personification}（拟人）：以人为喻（\textit{The wind whispered.}）。文学术语再配一对：\textbf{dead metaphor}（死喻——早已意识不到是比喻的比喻：\textit{the foot of the mountain、the arm of a chair}）。本讲义多处谈“词义引申的认知隐喻”（如 bitter 的味觉通感），它们大多是活着的 metaphor。
\end{gram}

\begin{origin}
希腊语 \textit{metaphora}“\textbf{转运、携带而过}”：\textit{meta-}（越过、转换）+ \textit{pherein}“\textbf{携带}”（原始印欧语 *\textit{bher-}“携带、生育”，日耳曼支即 \textbf{bear}“承担；生育”，拉丁支 \textit{ferre}“带来”$\to$\textbf{transfer, differ, offer, prefer, suffer}）。绝妙的对应：拉丁语用\textbf{同样的构词}造出 \textbf{transfer}（\textit{trans-}$+$\textit{ferre}“运过去”）——\textbf{metaphor 与 transfer 是希拉两种语言各自造出的同一个词}，“隐喻”就是意义的“转运”。同根 \textit{pherein} 还给了 \textbf{periphery}（peri- 周围携带$\to$周边）、\textbf{phosphorus}（携带光者$\to$磷）。
\end{origin}

\vw{mild}{/maɪld/}{adj.}

\sense{adj.}{不浓的，淡味的；暖和的；温和的；不严重的}
\col{a mild climate / winter（暖和）}
\ex{The city is famous for its mild, wet winters.}
\col{mild cheese / flavor / sauce（清淡）}
\ex{The curry is mild enough for children.}
\col{a mild manner / criticism（温和）}
\ex{He was a mild man who never raised his voice.}
\col{mild discomfort / surprise（轻微的）}
\ex{She felt mild discomfort in her knee after the run.}

\begin{gram}{mild 及其反义阵营的“分维对立”}
mild 是\textbf{多维弱化词}，它的反义词随维度更换：味道 $\leftrightarrow$ \textbf{strong / spicy}；气候 $\leftrightarrow$ \textbf{severe / harsh}（严寒）；性格 $\leftrightarrow$ \textbf{fierce / tough}（凶悍）；程度 $\leftrightarrow$ \textbf{serious / severe}（严重）。这种“一形容词多反义词”是词汇深度学习的标志，也是七选五阅读中利用反义衔接解题的钥匙（\textit{mild\dots however, harsh\dots}）。
\end{gram}

\begin{origin}
古英语 \textit{milde}“温柔的”，出自原始日耳曼语 *\textit{milthja-}，最终可溯至原始印欧语 *\textit{mel-}“\textbf{软}”——它在拉丁一支的伟大后裔是 \textit{mollis}“柔软的”：\textbf{mollify}（使软$\to$安抚）、\textbf{emollient}（润肤软化剂）、\textbf{mollusk}（软体动物）；希腊一支给出 \textit{maldus}? 否——但立陶宛语 \textit{minkštas}“软”与之呼应。日耳曼本支另有 \textbf{melt}（融化“使软”）疑为近亲。“温和”在词源上就是“\textbf{软}”：气候不硬、性格不硬、味道不硬——一个“软”字，四面开花。
\end{origin}

\vw{mixture}{/ˈmɪkstʃər/}{n.}

\sense{n.}{混合物；混合，结合}
\col{a mixture of A and B / a mixture of emotions}
\ex{The paste is a mixture of flour, eggs and milk.}
\ex{She listened with a mixture of amusement and alarm.}
\col{mix / mixed（动词／形容词）}
\ex{Mix the flour with a little water.}
\ex{The concert drew a mixed crowd of young and old.}

\begin{gram}{misc-“混合”词族}
拉丁语 \textit{miscere}“混合”（原始印欧语 *\textit{meik-}）在英语中藏得很深：\textbf{mixture}、\textbf{miscellaneous}（混杂的“杂项的”：\textit{miscellaneous expenses} 杂费）、\textbf{miscible}（可混溶的，化学术语）、\textbf{promiscuous}（pro- 向前乱混$\to$滥交的； indiscriminate）、以及意想不到的 \textbf{meddle}（插手“搅和进去”）与法语借词 \textbf{medley}（混泳接力；大杂烩）。凡见 \textit{misc-/medl-}，先想“搅在一起”。
\end{gram}

\vw{mood}{/muːd/}{n.}

\sense{n.}{情绪，心情；气氛，氛围}
\col{in a good / bad mood（心情好／坏）}
\ex{Careful --- the boss is in a terrible mood this morning.}
\col{be in the mood for sth / to do sth（有心情做）}
\ex{I'm not in the mood for arguing; let's talk tomorrow.}
\col{spoil / lighten the mood（破坏／活跃气氛）}
\ex{Her joke lightened the mood at the tense meeting.}
\col{the mood of the country / the times（民意／时代氛围）}
\ex{The novel captures the mood of America in the 1930s.}

\begin{gram}{mood 的语法学兄弟：语气}
\textbf{mood} 在语言学里专指“语气”：indicative mood（陈述语气）、imperative mood（祈使）、subjunctive mood（虚拟语气）。为什么“心情”与“语气”同词？因为语法的 mood 表达\textbf{说话人对所言的心理态度}（事实？愿望？命令？）——同一词根的两层开花。顺带正名：本讲义各 Language Focus 所谈“虚拟语气”，即 subjunctive mood，词源上就是“假设心情”的语法化。
\end{gram}

\begin{origin}
古英语 \textit{mōd}“\textbf{心、心灵、勇气}”（原始日耳曼语 *\textit{mōdaz}；德语 \textbf{Mut}“勇气”——\textit{gemütlich}“心情舒畅的”即其形容词；古诺斯语同词进入英语后还留下 \textbf{moody}“喜怒无常的”“心绪主导的”）。中古英语里 mōd 泛指“内心状态”，后窄化为“一时情绪”。它的复合词曾极其兴旺：\textbf{goodbye} $=$ God be with ye 的磨损形，而 \textit{mood} 退出复合词后，其位置由 \textbf{mind}（古英语 \textit{gemynd}“记忆、心”）接管。德语至今说 \textit{Hochmut}（高心$\to$傲慢）、\textit{Großmut}（大心$\to$高尚），英语的 mood 只剩下心里的一小片天气。
\end{origin}

\vw{multiple}{/ˈmʌltɪpl/}{adj.}

\sense{adj.}{数量多的，多样的}
\col{multiple choices / causes / injuries}
\ex{The accident had multiple causes, human and mechanical.}
\ex{He sustained multiple injuries in the fall.}
\col{multiply（动词：乘；成倍增加）}
\ex{Multiply 7 by 6 and you get 42.}
\ex{Our worries multiplied as the deadline approached.}

\begin{gram}{“多”的量级谱与 -plex 词族}
\textbf{multiple} $=$ \textit{multi-}（多）$+$ \textit{-plex}（\textbf{折叠}！即 complex 词条的 \textit{plec-}“编织”词根）——“多重折叠的”。由此解开一组数学词：\textbf{simple}（一折）、\textbf{double / triple / quadruple}（两折／三折／四折）、\textbf{multiple}（多折）、\textbf{multiply}（使成多折$\to$乘）、\textbf{multiplication}（乘法）。而“众多”形容词的量级：\textbf{several}（几个）$<$ \textbf{multiple / many}（多个）$\le$ \textbf{numerous}（许多，书面）$<$ \textbf{countless / innumerable}（无数）。
\end{gram}

%================================
\section{N}
%================================

\vw{native}{/ˈneɪtɪv/}{n.\ \&\ adj.}

\sense{n.}{出生于某国（或某地）的人；本地人}
\col{a native of sp（某地人）}
\ex{He is a native of Chengdu but works in Shenzhen.}
\col{a native speaker of English（母语者，高频）}
\ex{Ask a native speaker to check the collocations.}

\sense{adj.}{出生地的；本地的；土著的；土产的}
\col{one's native land / language / country}
\ex{She returned to her native village after forty years abroad.}
\col{native plants / species（土生物种）}
\ex{The panda is native to central China.（be native to 固定搭配）}
\col{go native（殖民时代旧语：入乡随俗，今含贬义慎用）}

\begin{gram}{nasc-“出生”词族与 gen- 词族的会师}
拉丁语 \textit{nasci}“出生”（其分词 \textit{natus}；实为 \textit{gnasci} 的 g 脱落形——与 gen-“生育”词根（见 generate 词条）\textbf{本是同一个词根}！）家族庞大：\textbf{native}（出生地的）、\textbf{nation}（同族出生者$\to$民族；国家）、\textbf{nature}（天生之性$\to$自然；本性）、\textbf{natural}、\textbf{natal}（出生的：\textit{prenatal} 产前）、\textbf{innate}（in- 内生$\to$天生的）、\textbf{cognate}（co- 同生$\to$同源的；同族词）、\textbf{nascent}（正出生的$\to$新兴的）、\textbf{Renaissance}（再生$\to$文艺复兴）。学与生在词源上本是一家。
\end{gram}

\vw{nourish}{/ˈnʌrɪʃ/}{vt.}

\sense{vt.}{培养，助长；抚养，滋养}
\col{nourish sth（希望、仇恨、才华等抽象物）}
\ex{For years she nourished the hope of returning to her hometown.}
\ex{Good teachers nourish curiosity rather than punish mistakes.}
\col{be nourished by（由……滋养）}
\ex{The valley is nourished by meltwater from the mountains.}
\col{nourishing / nourishment（滋养的／营养）}
\ex{A nourishing breakfast helps children concentrate.}
\ex{The soil provides constant nourishment for the vines.}

\begin{origin}
经由古法语 \textit{nourir} $\leftarrow$ 拉丁语 \textit{nutrire}“\textbf{哺乳、喂养}”（词源或与 \textit{nūtus}“点头示意”无关，一说出自拟声 \textit{nu-}“低哼哄婴”）。它的一族皆是“哺育”：\textbf{nourish}、\textbf{nurse}（乳母$\to$护士；动词护理；\textit{nursing a grudge} 怀恨“哺乳旧怨”！）、\textbf{nurture}（养育；培育）、\textbf{nutrition / nutritious}（营养／有营养的）、\textbf{nutrient}（养分）、\textbf{nurture} 的科学对词 \textit{nature vs.\ nurture}（先天对后天）——这对心理学名言正建立在“出生”与“哺乳”两个词根的对照上（nature 见 native 词条）。
\end{origin}

\vw{numerous}{/ˈnjuːmərəs/}{adj.}

\sense{adj.}{众多的，许多的（书面语）}
\col{numerous occasions / attempts / studies}
\ex{Numerous studies have confirmed the link between sleep and memory.}
\ex{She has been to that island on numerous occasions.}
\col{词族：number / numeral / innumerable / enumerate}
\ex{The stars are innumerable on a clear mountain night.（无数的）}
\ex{Please enumerate the reasons for your objection.（逐一列举）}

\begin{origin}
拉丁语 \textit{numerus}“\textbf{数、数目}”（词源成谜，或与 \textit{nem-}“分配、取份”相关——数目即“分配的单位”？此说亦解希腊 \textit{nemein}“分配”：\textbf{economy}$<$\textit{oikonomia}“家业分配”、\textbf{astronomy}?（否，星之“法则” \textit{nomos} 出于同根“分派秩序”！竟也相通））。\textit{numerus} 家族：\textbf{number}、\textbf{numeral}（数字符号）、\textbf{numerous}、\textbf{enumerate}（e- 数出$\to$列举）、\textbf{innumerable}（数不清的）、\textbf{supernumerary}（定额之外的$\to$临时演员；冗员）。罗马计数法无零、无位值，\textit{numerus} 的“份”义遂与日耳曼 \textbf{nimble}? 无关——数字词的族谱远比算术复杂。
\end{origin}

%================================
\section{O}
%================================

\vp{as opposed to}

\sense{prep. 短语}{（表示对比）而，相对于}
\ex{I'd prefer to arrive early, as opposed to rushing at the last minute.}
\ex{The company sells software, as opposed to hardware.}
\ex{The tax hits spending, as opposed to income.}

\begin{gram}{as opposed to / rather than / instead of}
三者都可表“而不是”，但语法性格不同：\textbf{as opposed to} 最正式，强调\textbf{对照}，后接名词／动名词；\textbf{rather than} 强调\textbf{取舍}，可接多种平行结构（\textit{He ran rather than walked.} 注意than后动词形式与前面平行）；\textbf{instead of} 口语化，介词性质（后接名词／动名词：\textit{instead of complaining}）。写作对比句模板：\textit{A, as opposed to B, does\dots}（A 与 B 相对，做……）。
\end{gram}

\begin{gram}{opposed 的两副面孔}
\textbf{opposed} 自身有两义，需借介词分辨：\textbf{be opposed to sth}“\textbf{强烈反对}”（表态度：\textit{Residents are opposed to the plan.}）；\textbf{as opposed to}“\textbf{截然不同、与之相对}”（表对照，无反对义）。一词两义皆出于 \textit{ob-}（逆）$+$ \textit{ponere}（放置）——“放在对面”：态度上针锋相对，或仅是位置上互为参照。
\end{gram}

\vp{out of this world}

\sense{adj. 短语}{好（或美等）得不得了，非凡，呱呱叫}
\ex{The dessert was out of this world --- I asked for the recipe.}
\ex{Their performance in the final was simply out of this world.}

\begin{gram}{out of 短语的“超离”修辞}
\textbf{out of this world} 字面“出了这个世界”，以夸张言其好。同构夸张家族：\textbf{over the moon}（欣喜若狂，义同“上月球”）、\textbf{on top of the world}（心花怒放）、\textbf{the best thing since sliced bread}（切片面包以来最好的，美俚）。反义夸张：\textbf{from hell}（\textit{the neighbors from hell} 恶魔邻居）、\textbf{a nightmare}（噩梦般的）。口语夸张词族是听力的“情绪标尺”，听到即判褒贬。
\end{gram}

\vw{outstanding}{/ˌaʊtˈstændɪŋ/}{adj.}

\sense{adj.}{优秀的，杰出的}
\col{an outstanding student / performance / contribution}
\ex{She was awarded a scholarship for her outstanding academic record.}
\col{an outstanding example（典范）}
\ex{The Great Wall is an outstanding example of ancient engineering.}

\sense{adj.}{突出的，明显的；（欠款、问题）未解决的（教材未列，商务高频）}
\col{an outstanding issue / debt / payment（悬而未决／未偿付的）}
\ex{Two points are still outstanding before the contract is signed.}
\ex{The company has outstanding debts of over a million yuan.}

\begin{gram}{“站出来”的词源逻辑与未偿付义}
\textbf{outstanding} $=$ out + standing（\textit{stā-}“站”词根，见 consist 词条）：“\textbf{站出来}”者即杰出者。“未偿付的”义出自会计意象：账目中\textbf{还立着没销}的款项（未清账仍“站”在簿上）。同构双义词还有 \textbf{striking}（突出的；打击的）与 \textbf{prominent}（突出的；凸出的）。一词多义多非偶然，而是\textbf{同一意象的不同投影}。
\end{gram}

\vw{overall}{/ˌəʊvərˈɔːl/}{adj.\ \&\ adv.}

\sense{adj.}{全面的，综合的}
\col{the overall picture / effect / cost}
\ex{The overall cost of the trip came to about eight thousand yuan.}
\ex{We must consider the overall effect of the policy, not one detail.}

\sense{adv.}{全部，总计；一般来说，大体上}
\ex{Overall, the concert was a great success despite the rain.}
\ex{The team spent, overall, about two million yuan on the project.}

\begin{gram}{overall / on the whole / in general / all in all}
四个“总而言之”的语体差：\textbf{overall} 最常用于数据总结与评论（可作定语\textbf{the overall winner} 总冠军）；\textbf{on the whole} 温和的“大体上”；\textbf{in general} 表常态规律（与 in particular 相对）；\textbf{all in all} 归结全篇的口语总结。图表作文的结论句首选 overall（雅思/考试通用）：\textit{Overall, sales rose in all three regions.}
\end{gram}

\vw{overcome}{/ˌəʊvəˈkʌm/}{vt.\ (overcame, overcome)}

\sense{vt.}{克服，解决；战胜}
\col{overcome difficulties / obstacles / a problem}
\ex{She overcame her fear of flying through slow practice.}
\col{overcome the enemy / one's rivals（战胜）}
\ex{The home team overcame a ten-point deficit to win.}

\sense{vt.}{受到……的极大影响}
\col{be overcome by / with emotion, grief, fumes（被情感／烟气压倒）}
\ex{Overcome with emotion, she could hardly finish the sentence.}
\ex{Several miners were overcome by toxic fumes.}

\begin{gram}{get over / overcome / conquer 的程度阶}
\textbf{get over}（口语：从疾病、失恋、震惊中\textbf{缓过来}：\textit{get over a cold / a breakup}）；\textbf{overcome}（正式：\textbf{克服}困难、恐惧）；\textbf{conquer}（最强：\textbf{攻克、征服}，常喻体：\textit{conquer Mount Qomolangma}；\textit{conquer one's addiction}）。注意 \textit{be overcome with} 的“被情感淹没”义\textbf{不可}换成 get over——方向相反：一个是压倒，一个是恢复。
\end{gram}

\vw{overthrow}{/ˌəʊvəˈθrəʊ/ vt.\ \ /ˈəʊvəθrəʊ/ n.}{vt.\ \&\ n.}

\sense{vt.}{推翻，打倒}
\col{overthrow a government / a regime}
\ex{The rebels attempted to overthrow the government in a single night.}
\ex{The old theory was overthrown by new evidence.（喻：推翻旧说）}

\sense{n.}{推翻，打倒}
\ex{The overthrow of the dynasty was plotted for years.}

\begin{gram}{over- 动词家族与名动重音（三度重逢）}
\textbf{overthrow} 属 over- 前缀动词家族：\textbf{overcome}、\textbf{overtake}（追上；超车；突然袭击）、\textbf{overwhelm}（倾覆；压倒）、\textbf{overlook}（俯视$\to$忽视）、\textbf{overhear}（无意听到）、\textbf{overrule}（否决）。名词化后重音前移：\textit{over\textbf{throw}}（动）$\to$\textit{\textbf{over}throw}（名）——与 advocate、estimate、digest 同规则，本讲义第三次也是最后一次提醒此构词法，务必内化。
\end{gram}

\vw{owe}{/əʊ/}{vt.}

\sense{vt.}{欠（债）；欠（情）}
\col{owe sb sth / owe sth to sb（欠某人某物）}
\ex{I owe you an apology for what I said yesterday.}
\ex{The young writer owes much to her first editor.（深受其惠）}
\col{owe it to sb to do sth（对某人有义务做）}
\ex{We owe it to our children to leave them a cleaner world.}

\sense{vt.}{归功于（owe sth to sth）}
\ex{She owes her success to years of quiet practice.}

\begin{gram}{owe / owing to / owe it to oneself}
\textbf{owing to} $=$ due to“由于”（\textit{Owing to the fog, all flights were delayed.}），与“欠”义已语法化为独立介词，多用于句首；而 \textbf{owe it to oneself to do}“理应为自己做”（\textit{You owe it to yourself to take a rest.}）。三处考点：（1）\textit{owe} 后接\textbf{双宾语}（owe sb sth）；（2）\textit{owe\dots to\dots} 表归功；（3）owing to / due to / because of 的介词性（后接名词）。
\end{gram}

\begin{origin}
古英语 \textit{agan}“\textbf{拥有}”（原始日耳曼语 *\textit{aganan}，与德语 \textbf{besitzen?}——其近亲为古高地德语 \textit{eigun}“所有物”$\to$\textbf{own}“自己的”）。语义漂移堪称奇观：“拥有”$\to$\textbf{“欠别人拥有之物”}$\to$“欠”；而“负有义务”的\textbf{过去式} \textit{ought}（\textit{ahte}“曾欠”$\to$“本应”）竟脱胎换骨成了\textbf{情态动词}——“应该”的本义是“\textbf{欠着}”！同族 \textbf{own}（拥有；承认“认作己有”）与之共享家谱。一个“拥有”的词根，长出“欠”“该”“自己的”三个方向，活脱一部 mini 经济思想史。
\end{origin}

%================================
\section{P}
%================================

\vw{patience}{/ˈpeɪʃns/}{n.}

\sense{n.}{耐心，忍耐力}
\col{have the patience to do sth}
\ex{It takes real patience to teach children to swim.}
\col{run out of / lose patience (with sb/sth)}
\ex{The customers were running out of patience with the slow service.}
\col{with patience（耐心地）；test sb's patience（考验耐性）}
\ex{The toddler tested her patience all afternoon.}
\col{patient（形容词：有耐心的；名词：病人）}
\ex{Be patient with beginners --- everyone starts somewhere.}
\ex{The doctor sees fifteen patients a day.（名词）}

\begin{gram}{patient 的“病人”与“耐心的”：一词两义的受苦内核}
\textbf{patient} 作名词“病人”与形容词“有耐心的”并非巧合撞形，而是同一词根的两义：拉丁语 \textit{pati}“\textbf{承受、受苦}”的现在分词 \textit{patiens}——“正在承受的人”既是\textbf{病人}（承受病痛），也是\textbf{有耐心者}（承受等待与挫折）。由此解开一批词：\textbf{passion}（本义“受难”——\textit{the Passion of Christ} 基督受难；后泛化为“激情”）、\textbf{passive}（承受的$\to$被动的）、\textbf{compassion}（com- 共同承受$\to$同情）、\textbf{compatible}（可共处承受$\to$兼容的）、\textbf{impatient}（不能承受$\to$不耐烦）。记牢“pati $=$ 受”，这一族全部贯通。
\end{gram}

\vw{perceive}{/pəˈsiːv/}{vt.}

\sense{vt.}{注意到，意识到}
\ex{We perceived a slight change in her tone after the phone call.}
\col{perceive sth as sth / perceive that\dots}
\ex{The change was widely perceived as a step backward.}

\sense{vt.}{将……理解为，将……视为}
\col{be perceived as（被视为，新闻高频被动）}
\ex{Nuclear power is often perceived as risky, whatever the data says.}
\col{perception（名词：感知；看法）}
\ex{The drug alters one's perception of time.}
\ex{Public perception of the industry has improved.（公众看法）}

\begin{gram}{ceive/cept“取”词族：一对成对出现的词干}
拉丁语 \textit{capere}“\textbf{抓取}”的动词干在法语借词中化作 \textbf{-ceive}，其过去分词化作 \textbf{-cept}，遂成英语中成对出现的大家庭：\textbf{receive/reception}（取回$\to$收到；接待）、\textbf{deceive/deception}（骗取$\to$欺骗）、\textbf{conceive/concept}（合取$\to$构想；概念）、\textbf{perceive/perception}（per- 彻底取$\to$感知）、\textbf{except/exception}（取出去$\to$除外）、\textbf{accept}（取向自己$\to$接受）、\textbf{capture/captive}（拉丁原形直接借入$\to$捕获；俘虏）。动词用 -ceive，名词形容词用 -cept/-ception，一条规则收一串词。
\end{gram}

\vw{perspective}{/pəˈspektɪv/}{n.}

\sense{n.}{透视法}
\ex{Renaissance painters mastered linear perspective to create depth.}

\sense{n.}{态度，观点；客观判断力}
\col{from sb's perspective / from the perspective of（从……的视角）}
\ex{From the students' perspective, the exam was unfair.}
\col{put sth in / into perspective（正确看待，摆正位置）}
\ex{Traveling puts your daily worries into perspective.}
\col{keep a sense of perspective（保持客观分寸）}
\ex{It's only a game --- keep a sense of perspective.}

\sense{n.}{景观，远景}
\ex{The tower offers a fine perspective of the whole valley.}

\begin{gram}{spec-“看”词族总动员（二）}
\textbf{perspective} $=$ \textit{per-}（透过）$+$ \textit{specere}“\textbf{看}”（见 insight 词条的“看”家谱）——“\textbf{透过去看}”：透视法如此，观点亦然（透过某个立足点看）。家族全员：\textbf{inspect}（往里看$\to$检查）、\textbf{respect}（反复看$\to$尊重）、\textbf{suspect}（由下看$\to$怀疑）、\textbf{prospect}（向前看$\to$前景）、\textbf{retrospect}（向后看$\to$回顾：\textit{in retrospect} 回想起来）、\textbf{spectator}（观看者）、\textbf{spectacle}（景观；眼镜）、\textbf{spectrum}（光谱“所见之范围”）。写作金句：\textit{Perspective is everything.} 立场即所见。
\end{gram}

\vw{pixel}{/ˈpɪksl/}{n.}

\sense{n.}{像素}
\col{pixels per inch（每英寸像素，屏幕参数）}
\ex{The camera shoots 48 million pixels per image.}
\col{a pixelated image（像素化模糊的）}
\ex{Zoom in too far and the photo turns pixelated.}

\begin{gram}{合成缩略词（blend）：两个词“对半拼”}
\textbf{pixel} $=$ \textbf{pic}ture（图片）$+$ el\textbf{ement}（元素）各取一半熔合而成——这类“熔合词”（blend / portmanteau）是现代科技造词的主力：\textbf{smog}（smoke + fog）、\textbf{brunch}（breakfast + lunch）、\textbf{motel}（motor + hotel）、\textbf{internet}（inter- + network）、\textbf{blog}（web + log）、\textbf{modem}（modulator + demodulator）。与之相对的是\textbf{首字母缩略}（acronym）：\textbf{laser, radar, NASA}。掌握这两种构词法，科技文生词的恐惧减半。
\end{gram}

\vw{plain}{/pleɪn/}{adj.\ \&\ n.}

\sense{adj.}{朴素的，简单的；清楚的}
\col{plain food / clothes（素净）}
\ex{The monks eat plain food and wear plain robes.}
\col{in plain English / words（明白说）}
\ex{Could you explain it in plain English, without the jargon?}
\col{plain honesty（直白）}
\ex{Her plain speaking sometimes offended people.}
\col{as plain as day（一清二楚）}
\ex{It was as plain as day that he was lying.}

\sense{n.}{平原}
\ex{Herds of wild horses roam the great plains.}

\begin{origin}
经由古法语 \textit{plain} $\leftarrow$ 拉丁语 \textit{planus}“\textbf{平的}”（原始印欧语 *\textit{pelə-}“平坦”，同根 \textbf{floor}?（日耳曼 \textit{flōr}“平地”是其远亲）、\textbf{plane}（刨子；飞机“平飞的翼面”）、\textbf{plan}（平面图$\to$计划）、\textbf{plate / plateau}（平板／高原）、\textbf{plain}（平原）与意外亲戚 \textbf{explain}（ex- 摊平$\to$把话摊开$\to$解释——\textbf{explain 与 plain 同根}！把道理\textbf{摊平}了自然“明白”）。“朴素”义取“平实无饰”，“清楚”义取“平铺直叙”。一词根，从高原摊到说明书。
\end{origin}

\vw{political}{/pəˈlɪtɪkl/}{adj.}

\sense{adj.}{政治的，政府的，政权的；政党的，党派的}
\col{a political party / system / crisis}
\ex{The country is passing through a deep political crisis.}
\col{political prisoners / asylum（政治犯／政治庇护）}
\ex{The dissident spent years as a political prisoner.}
\col{a political animal（政治动物，好权术者）}
\ex{He admits he is a political animal who thrives on debate.}
\col{politics / politician / polity（词族）}
\ex{She studied politics before entering parliament.}
\ex{Career politicians often begin as local councilors.}

\begin{origin}
经由拉丁语 \textit{politicus} $\leftarrow$ 希腊语 \textit{politikos}“\textbf{城邦公民的}”，出自 \textit{polis}“\textbf{城邦}”——古希腊的“政治”本义就是“\textbf{城邦里众人之事}”，亚里士多德名言“人是城邦的动物”（\textit{zōon politikon}）即用此词。\textit{polis} 家族在英语中灯火通明：\textbf{police / policy}（城邦治理的两只手：治安$\to$警察；治理之策$\to$政策）、\textbf{politician}（从政者）、\textbf{cosmopolitan}（\textit{kosmos} 世界 + \textit{polis} 城邦$\to$世界公民）、\textbf{metropolis}（\textit{mētēr} 母 + \textit{polis}$\to$母城、大都会）、\textbf{megalopolis}（巨大都市带）、地名后缀 \textbf{-polis}（Annapolis, Minneapolis）与 \textbf{-ville}?（法语系，勿混）。识得 polis，城市与政治一通百通。
\end{origin}

\vw{possess}{/pəˈzes/}{vt.}

\sense{vt.}{具有（特质）；拥有；支配，控制}
\col{possess sth（品质、能力）}
\ex{She possesses a rare gift for putting people at ease.}
\col{be possessed of（书面：具有）／ be possessed by（被……支配、附体）}
\ex{He was possessed of enormous energy and ambition.}
\ex{In folk tales, a person could be possessed by an evil spirit.}
\col{possess sb to do?（美口语：What possessed you to\dots 你着了什么魔）}
\ex{What possessed you to dye your hair green?}
\col{possession（名词：拥有；财产）}
\ex{The island was once in the possession of Portugal.（in possession of 拥有）}

\begin{gram}{possess 与 own 的语法性格差}
\textbf{possess} 正式，多接\textbf{抽象品质}（qualities, courage, evidence），可用于被动（\textit{be possessed of}）；\textbf{own} 日常，多接\textbf{具体财物}，不用被动（\ding{55} \textit{A car is owned by me} 生硬），其形容词 \textbf{own}（自己的）另成体系（\textit{my own house}）。固定区分：\textit{in possession of}（主动拥有）vs.\ \textit{in the possession of}（\textbf{被}他人拥有，多 the）——冠词之差，主客易位，是精密阅读的经典陷阱。
\end{gram}

\begin{origin}
拉丁语 \textit{possidēre}“\textbf{坐拥}”：\textit{pos-}（后期拉丁语 $=$ \textit{potis}“能够；主人”，同根 \textbf{possible}、\textbf{potent}、\textbf{despot}“主人$\to$暴君”）$+$ \textit{sedēre}“\textbf{坐}”（\textit{sed-}“坐”词根：\textbf{sediment}（坐沉之物$\to$沉积物）、\textbf{sedentary}（久坐的）、\textbf{sedate}（安坐的$\to$镇静的；肃静的）、\textbf{preside}（坐在前面$\to$主持）、\textbf{reside}（坐回$\to$居住；归属）、\textbf{subsidy}（坐在下面支助$\to$补贴））。“拥有”的本义是\textbf{以主人的身份坐镇其上}——罗马法的产权观（占有即坐着）尽在一词。而 \textbf{obsess}（ob- 坐在对面$\to$围坐施咒$\to$使着迷）是它的镜像兄弟。
\end{origin}

\vw{precise}{/prɪˈsaɪs/}{adj.}

\sense{adj.}{细致的，精细的；准确的，精确的；恰好的}
\col{precise measurements / calculations（精确）}
\ex{The bridge was built to precise measurements.}
\col{precise instructions / details（细致）}
\ex{Follow the precise instructions on the packet.}
\col{to be precise（确切地说，插入语）}
\ex{It happened in 1975, to be precise, on the night of July the twelfth.}
\col{precision（名词：精度；精密）}
\ex{The parts are machined with great precision.}

\begin{gram}{precise / accurate / exact 的微差}
\textbf{accurate}“准确”（与事实相符，无差错：\textit{an accurate translation}）；\textbf{exact}“确切”（丝毫不差，常重数量细节：\textit{the exact number}）；\textbf{precise}“精确”（精细到高分辨率，重方法与刻度：\textit{precise instruments}）。口语中三词常互换，考试中依搭配锁义：\textit{precise} 配 instruments/measurements，\textit{accurate} 配 information/data，\textit{exact} 配 number/time/words。插入语 \textit{to be precise} 与 \textit{to be exact / accurate} 相比更重“给出更细刻度”。
\end{gram}

\vw{prediction}{/prɪˈdɪkʃn/}{n.}

\sense{n.}{预言，预测}
\col{make a prediction (about sth)}
\ex{It is too early to make a prediction about the election.}
\col{predictions come true / prove accurate（预言成真／应验）}
\ex{Few predictions about technology come true on schedule.}
\col{predictable / unpredictable（可预测的／变幻莫测的）}
\ex{The endings of his films are never predictable.}
\col{predict（动词）}
\ex{No one could predict how the war would end.}

\begin{gram}{pre-“预先”前缀小集}
\textbf{predict} $=$ \textit{pre-}（预先）$+$ \textit{dicere}“说”（\textit{dic-}“言”词根，见 contradictory 词条：“预先说”即预言）。pre- 家族速览：\textbf{prevent}（预先来$\to$阻止）、\textbf{prepare}（预先备置）、\textbf{preview}（预览）、\textbf{previous}（在前的）。注意 \textit{pre-}“事先”与 \textit{pro-}“向前”方向不同：predict（事先说）对 proceed（向前走）。另外 \textit{predict} 与 \textit{forecast}（fore + cast“往前掷”，见 cast 词条）是拉丁、日耳曼各造的“预言”对词，正如 metaphor/transfer 之对——英语常给同一概念配希拉、日耳曼双份表达，一雅一俗。
\end{gram}

\vw{property}{/ˈprɒpəti/}{n.}

\sense{n.}{性质，特性}
\col{a chemical / physical / healing property}
\ex{Scientists studied the magnetic properties of the new material.}
\ex{Many plants have medicinal properties.}

\sense{n.}{所有物，财产；不动产}
\col{private / public property}
\ex{Respect other people's private property.}
\col{intellectual property（知识产权）；property prices（房价）}
\ex{The firm specializes in intellectual property law.}
\ex{Property prices in the city have doubled in a decade.}

\begin{gram}{property 的“属性”与“财产”：一词两义的哲学}
\textbf{property} 的两义共享词根 \textit{proprius}“\textbf{自己特有的}”（见 appropriate 词条）：属于某物\textbf{特有}的品质是“属性”，属于某人\textbf{特有}的物是“财产”。科技英语中 property 几乎只作“性质”复数（\textit{the properties of water}）；法律经济语境中多指“财产、不动产”（英式常特指\textbf{房产}）。写作辨义技巧：看谓语——\textit{The property belongs to\dots} 财产；\textit{The property of X is\dots} 属性。
\end{gram}

%================================
\section{R}
%================================

\vw{rare}{/reər/}{adj.}

\sense{adj.}{稀少的，罕见的；稀罕的，珍贵的}
\col{a rare event / species / occurrence}
\ex{Snow is a rare event in this southern city.}
\ex{The wetlands shelter several rare species of crane.}
\col{it is rare for sb to do sth（很少见某人做）}
\ex{It is rare for him to arrive late, let alone miss class.}
\col{rarely（副词：很少——句首倒装!）}
\ex{\textbf{Rarely have I} seen such a courageous young player.}
\col{a rare steak（嫩牛排，教材未列，餐饮必知）}
\ex{I'd like my steak rare, please.}

\begin{gram}{rarely 句首倒装：否定副词的仪式}
\textbf{rarely / seldom / never / little / hardly / scarcely / no sooner} 置于句首时，句子须\textbf{部分倒装}（助动词提前）：\textit{Rarely \textbf{does he} complain.} / \textit{Not once \textbf{did she} look back.} / \textit{Hardly \textbf{had we} started when it began to rain.}（hardly\dots when 固定搭配）。这与 only 状语前置（\textit{Only then \textbf{did I} realize}）同属一个倒装家族——语法填空与改错的高频考点，倒装的本质是\textbf{把否定拉到句首聚焦}。
\end{gram}

\vw{raw}{/rɔː/}{adj.}

\sense{adj.}{未经加工的，自然状态的；生的，未烹制的}
\col{raw materials（原材料，工业高频）}
\ex{The country imports most of its raw materials.}
\col{raw fish / vegetables（生的）}
\ex{Sushi is made with raw fish of the highest freshness.}
\col{raw data / footage（原始数据／影像）}
\ex{The team analyzed the raw data from ten years of surveys.}
\col{raw talent（未经雕琢的天赋）；a raw deal（不公平待遇，习语）}
\ex{The young singer has raw talent but needs training.}
\ex{Laid-off workers got a raw deal from the management.}

\begin{origin}
古英语 \textit{hrēaw}“\textbf{生的、血淋淋的}”（原始日耳曼语 *\textit{hrawaz}，与德语 \textbf{roh} 同源），最终出自原始印欧语 *\textit{krewə-}“\textbf{生肉、血肉}”——一个惊人同根词是拉丁语 \textit{crūdus}“生的”：其后代 \textbf{crude}（原油 \textit{crude oil}“生的油”！ crude 与 raw 是希拉两种语言的\textbf{同义词对}）、\textbf{cruel}（血淋淋的$\to$\textbf{残酷}——cruel 的词源竟是“生肉”！）、\textbf{crudeness /crudity}。希腊同根 \textit{kreas}“肉”藏在 \textbf{pancreas}（pan + kreas“全肉腺”$\to$胰腺）里。生的、残忍的、胰腺，同一块史前生肉。
\end{origin}

\vw{realistic}{/ˌriːəˈlɪstɪk; ˌrɪəˈlɪstɪk/}{adj.\ \&\ reality /riˈæləti/ n.}

\sense{adj.}{逼真的，栩栩如生的；现实的，实际的；明智的}
\col{a realistic painting / model（逼真）}
\ex{The battle scenes are so realistic that some viewers looked away.}
\col{be realistic about sth（面对现实）}
\ex{You must be realistic about your chances of winning.}
\col{a realistic goal / target / expectation（切实可行的）}
\ex{Finishing in three months is a realistic target.}
\col{realize（实现；意识到）}
\ex{She finally realized her dream of opening a bookshop.}

\sense{reality n.}{现实，实际情况；事实，实际经历}
\col{face reality / in reality（实际上）}
\ex{In reality, the project cost twice the original estimate.}
\col{escape from reality（逃避现实）}
\ex{Gaming became his way of escaping from reality.}

\begin{origin}
拉丁语 \textit{realis}“\textbf{物的}”，出自 \textit{res}“\textbf{物、事}”（词源晦暗，或与 \textit{rē-}“物、财”相关——罗马法术语 \textit{res} 遍地开花：\textbf{republic}（\textit{res publica} 公共之物$\to$共和国）、\textbf{real estate}（地产“实物产业”）、\textbf{realize}（使成物$\to$\textbf{实现}；认知上“见其为实”$\to$\textbf{意识到}）、\textbf{realism}（实在论；现实主义）、\textbf{realty / realtor}（房地产业）。中世纪经院哲学的\textbf{唯实论}（Realism，“共相为 real”）与唯名论之争使 real 成为哲学关键词；日常语的 realistic 由此带着“按事物本相”的基因。
\end{origin}

\vp{rely on}

\sense{v. 短语}{依赖，依靠；信任}
\col{rely on/upon sb/sth (to do sth)}
\ex{The villagers rely on the river for water and fish.}
\ex{You can rely on David to keep his word.}
\col{rely on sb/sth for sth（靠……获得）}
\ex{The region relies heavily on tourism for jobs.}
\col{reliable / reliance（词族）}
\ex{She is the most reliable member of the team.}
\ex{The country's reliance on imported oil has fallen.（reliance on）}

\begin{origin}
经由古法语 \textit{relier}“\textbf{再缚、系牢}”$\leftarrow$ 晚期拉丁语 \textit{religāre}（\textit{re-} 再 + \textit{ligāre}“\textbf{捆绑}”）：“依靠”即“\textbf{把自己再拴上去}”。\textit{ligāre}（原始印欧语 *\textit{leig-}“捆”）的家族皆与“绑”有关：\textbf{ally}（结盟“绑在一起”）、\textbf{league}（联盟；里格“一拃之距”？后者另有来源）、\textbf{league} 同源 \textbf{colleague}（col- 共同绑$\to$同事）、\textbf{obligation}（绑住$\to$义务）、\textbf{ligament}（韧带）、\textbf{ligature}（结扎线）。最耐人寻味的亲戚是 \textbf{religion}（宗教）：一词之解千载聚讼，西塞罗解为 \textit{relegere}“反复敬读”，拉克坦提乌斯与奥古斯丁解为 \textit{religare}“\textbf{重新系于人神}”——若取后说，\textbf{rely 与 religion 同根}：信仰与信赖，都是把自己交绑给什么。
\end{origin}

\vw{remarkable}{/rɪˈmɑːkəbl/}{adj.}

\sense{adj.}{非凡的，奇异的，显著的，引人注目的}
\col{a remarkable achievement / similarity / change}
\ex{The two portraits show a remarkable similarity.}
\ex{It is remarkable that no one was hurt in the accident.}
\col{remarkably（显著地；非凡地）}
\ex{The economy recovered remarkably quickly.}
\col{remark（动词：评论；名词：言论）/ It is worth remarking that\dots}
\ex{She remarked that the room had grown strangely quiet.}

\begin{gram}{-able / -ible 形容词的“值得”义与动词回溯}
\textbf{remarkable} $=$ remark（评论）$+$ -able“\textbf{可、值得}”——“值得评论的”$\to$“值得注意的”。同型推理：\textbf{respectable}（可敬的）、\textbf{admirable}（令人钦佩的）、\textbf{noticeable}（值得注意的）、\textbf{memorable}（值得记住的）、\textbf{preferable}（更可取的）。-able 后缀还保留了\textbf{及物关系}（remarkable things $=$ things that deserve remark），而被动义（\textit{a remarkable fact observed by all}）混入其中，这正是 -able/-ible 双形并存的古因之一。
\end{gram}

\vw{remote}{/rɪˈməʊt/}{adj.}

\sense{adj.}{偏远的，偏僻的；遥远的，久远的；远亲的}
\col{a remote village / area（偏远）}
\ex{Medical teams travel to remote villages twice a year.}
\col{in the remote past / future（久远）}
\ex{In the remote future, the sun will expand into a red giant.}
\col{a remote ancestor / relative（远亲）}
\ex{A remote ancestor of hers sailed to America in 1620.}
\col{a remote chance / possibility（可能性极小的，引申高频）}
\ex{There is a remote chance that the weather will clear up.}
\col{remote control（遥控）}
\ex{Where did you put the TV remote control?}

\begin{gram}{“远”的近义阶梯与 remote chance}
空间之远：\textbf{far}（普通）$\to$ \textbf{distant}（书面，见 D 部分）$\to$ \textbf{remote}（远而隔绝、难达）。概率之“远”是新托福与考研阅读的宠儿：\textit{a remote possibility / chance}（微乎其微的可能），反义 \textit{a strong / good likelihood}。注意 remote 的这层“小概率”义与空间无关，不可译“遥远的可能性”（\ding{55} 汉语搭配错误）。
\end{gram}

\vw{representative}{/ˌreprɪˈzentətɪv/}{adj.\ \&\ n.}

\sense{adj.}{典型的，有代表性的}
\col{be representative of（典型于，固定介词 of）}
\ex{This survey is representative of the views of young voters.}
\ex{Is this sample representative of the whole batch?}

\sense{n.}{典型人物，代表性人物；代表}
\col{a representative of / to sth（……的代表；驻……代表）}
\ex{She is our sales representative in the south of China.}
\ex{The union representatives met the management for talks.}
\col{a representative of the people（人民的代言人）}
\ex{Legislators see themselves as representatives of the people.}

\begin{gram}{represent 三义与 congressional vocabulary}
动词 \textbf{represent} 的三个面向：（1）\textbf{代表}（人）：\textit{The lawyer represented the company in court.}；（2）\textbf{表现、描绘}：\textit{The map represents the island as it was in 1800.}；（3）\textbf{意味着、相当于}：\textit{This contract represents months of negotiation.} 形容词 representative 表“典型的”，名词表“代表、议员”（美语 \textit{Representative} 众议员，\textit{House of Representatives} 众议院）。注意 representative 作形容词时重音与作名词相同，不属“名动异读”词对。
\end{gram}

\vw{reputation}{/ˌrepjuˈteɪʃn/}{n.}

\sense{n.}{名誉，名声}
\col{have a good / bad reputation for sth（因……而闻名）}
\ex{The restaurant has an excellent reputation for fresh seafood.}
\col{live up to one's reputation（名副其实）}
\ex{The hotel lived up to its reputation for luxury.}
\col{ruin / damage / enhance one's reputation}
\ex{One scandal nearly ruined his reputation.}
\col{reputation as（作为……的名声）}
\ex{She earned a reputation as a tough but fair negotiator.}

\begin{gram}{reputation / fame / prestige}
\textbf{fame} 中性之“名”（可好可坏？实多指知名度：\textit{fame and fortune}）；\textbf{reputation} 是他人口中的\textbf{评价}，有好坏之分（\textit{a reputation for honesty}）；\textbf{prestige} 崇高的\textbf{威望}（法语借词，原义“魔术戏法”！—— Prestige 一词的词源提醒我们：威望本是幻术）。短语 \textbf{of repute}（有名望的：\textit{a writer of repute}）与其反义 \textbf{disrepute} 同根。
\end{gram}

\begin{origin}
拉丁语 \textit{reputāre}“\textbf{再算、反复掂量}”：\textit{re-}（再）+ \textit{putāre}“\textbf{算、清理、想}”（原始印欧语 *\textit{pau-}“切、砍”——本义“修剪”，引申“清点、思量”）。\textit{putāre} 是个被低估的词根大户：\textbf{compute}（合算$\to$计算；computer）、\textbf{calculate?}（否，calculus 是石子）、\textbf{dispute}（dis- 分着算$\to$各执一词$\to$\textbf{争论}）、\textbf{repute}（再算$\to$掂量出的\textbf{名声}）、\textbf{deputy}（de- 代理算$\to$副手；deputy 的 -put 与 repute 同根！）、\textbf{amputate}（am- 环切$\to$截肢）、\textbf{putative}（推算的$\to$推定的）。“名声”在词源里是一笔\textbf{众人心中的反复核算}。
\end{origin}

\vw{reward}{/rɪˈwɔːd/}{vt.\ \&\ n.}

\sense{vt.}{奖励，奖赏}
\col{reward sb with sth / reward sb for sth}
\ex{The school rewarded her with a scholarship for her research.}
\ex{He was rewarded for his years of quiet service.}

\sense{n.}{奖励，回报；悬赏金}
\col{a reward for sth（注意 for：作为……的报酬／悬赏）}
\ex{The reward for information leading to an arrest rose to 50,000 yuan.}
\col{claim / offer a reward（认领／悬赏）}
\ex{The police offered a reward for the missing painting.}
\col{no reward without effort}
\ex{Success is the reward of hard work.}

\begin{gram}{reward / award：一字母之争}
\textbf{reward}“报偿、酬谢”——出于\textbf{应得}（行为、功劳之后的回报，可好可坏：\textit{reap the rewards} 收获回报）；\textbf{award}“授予、颁发”——出于\textbf{评判}（评奖授予：\textit{award sb a prize / be awarded the Nobel Prize}；名词“奖项”：\textit{an award ceremony} 颁奖典礼）。搭配锁定：\textit{Nobel Prize} 搭 award 不搭 reward；\textit{a reward for finding the lost cat} 用 reward。高考完形的一字母之差。
\end{gram}

\begin{origin}
经由古法语 \textit{rewarder}“\textbf{回报、注视}”的后代（\textit{rewardier}），本是 \textbf{regarder}“看顾”（re- 反复 + \textit{garder}“看守”，日耳曼词根 *\textit{wardōn}“守望”）的变体——与 \textbf{regard}（注视；尊敬）、\textbf{guard}（卫兵）、\textbf{ward}（病房；监护；\textit{warden} 看守人）、\textbf{warrant}／\textbf{guarantee}（见 G 部分，同根的“郑重看管之诺”）一家。\textbf{reward 的词源是“看顾”}：奖赏即“把目光（与财物）\textbf{回赠}给应得之人”；\textbf{regard} 则从“看”升华为“敬重”。看与予、看与敬，在这家词族里是一件事的两面。
\end{origin}

\vw{rhyme}{/raɪm/}{n.\ \&\ vt.\&vi.}

\sense{n.}{押韵词；押韵}
\col{a rhyme for sth（……的韵脚）}
\ex{``Day'' is a common rhyme for ``way''.}
\col{nursery rhymes（童谣）}
\ex{She sang nursery rhymes to put the baby to sleep.}

\sense{vi. \& vt.}{使押韵；和……同韵}
\ex{``Though'' rhymes with ``low''.（rhyme with 与……押韵）}
\ex{The poet rhymed ``moon'' with ``June''.}

\begin{gram}{rhyme / rhythm：拼写与概念的双重辨析}
\textbf{rhyme}“押韵”——\textbf{音的相同}（末尾音相同或相近：\textit{way / day}）；\textbf{rhythm}“节奏”——\textbf{音的时值安排}（轻重缓急的型式）。中文速记：\textbf{押韵管“同”，节奏管“序”}。诗学术语配套：\textit{alliteration}（头韵：\textit{safe and sound}）、\textit{assonance}（元音谐）、\textit{blank verse}（素体诗——有节奏无韵脚）。流行歌词表记 \textit{rhyme scheme}（韵式 abab / aabb）是歌词赏析的抓手。
\end{gram}

\begin{origin}
拼写骗局：中古英语本作 \textit{rime}（借自古法语 \textit{rime}，日耳曼源 \textit{rīm}“数目、系列”），17 世纪学者\textbf{误以为}它源自希腊语 \textit{rhythmos}（节奏），遂照希腊拼法添上 h 与 y——\textbf{rhyme 的 h 是一场误会}，是“民俗词源”（folk etymology，见 F 部分 folk 词条）在正字法上的胜利。真正的 \textbf{rhythm} $\leftarrow$ 希腊 \textit{rhythmos}“\textbf{测得的流动}”（\textit{rhein}“流”$+$表示“度量”的后缀，原始印欧语 *\textit{srew-}“流”：同根 \textbf{stream}?（日耳曼近亲）、\textbf{rheostat}（变阻器）、医学 \textbf{diarrhea}（dia- 穿流而过$\to$腹泻）、\textbf{rhea}（美洲鸵“奔流者”？实另源）。两个“音的”词，一个是被误认的日耳曼词，一个是真正的希腊流。
\end{origin}

\vw{rhythm}{/ˈrɪðəm/}{n.}

\sense{n.}{节奏，韵律，律动；规律}
\col{the rhythm of the music / drums（音乐节奏）}
\ex{The dancers moved to the rhythm of the drums.}
\col{the rhythm of life（生活的律动）}
\ex{Farm life follows the rhythm of the seasons.}
\col{circadian rhythm（生理节律，科普高频）}
\ex{Late nights disturb the body's circadian rhythm.}

\begin{gram}{rhythm 的语篇高频搭配}
科技与健康类阅读的常客：\textbf{sinus rhythm}（窦性心律，心脏语境）、\textbf{rhythm and blues}（节奏布鲁斯，R\&B——流行音乐体裁名）、\textbf{sense of rhythm}（节奏感：\textit{a natural sense of rhythm}）。形容词 \textbf{rhythmic / rhythmical}（有节奏的；前者更常用）：\textit{the rhythmic sound of the waves}。表“规律”时近义 \textbf{pattern / cycle}（见 cycle 词条），但 rhythm 强调\textbf{时间中的强弱起伏}，pattern 泛指模式。
\end{gram}

\vw{rigid}{/ˈrɪdʒɪd/}{adj.}

\sense{adj.}{死板的，僵硬的；固执的，僵化的}
\col{rigid rules / controls / attitudes（死板）}
\ex{The school's rigid rules left no room for creativity.}
\col{a rigid structure / material（刚性，物理）}
\ex{The frame must be rigid enough to hold the glass panels.}
\col{rigidity（名词：僵硬；僵化）}
\ex{The rigidity of the old system slowed every reform.}

\begin{gram}{rigid / stiff / tough / strict}
\textbf{rigid}：物体“刚性的不弯”；制度“僵死不变”（贬）；\textbf{stiff}：身体、关节、面料的“发僵”（\textit{stiff neck}），也指拘谨（\textit{a stiff smile}）；\textbf{tough}：坚韧抗造（\textit{tough material}；\textit{a tough decision} 艰难）；\textbf{strict}：专对\textbf{规则与人}的严格（\textit{strict parents / rules}）。搭配避坑：“严格的规定”可说 rigid rules 或 strict rules，但“严格的老师”只用 strict。
\end{gram}

\begin{origin}
拉丁语 \textit{rigidus}“\textbf{僵硬的}”，出自动词 \textit{rigēre}“僵直、发抖？”（实义“僵冷”——与 \textit{frīgēre}“冷”在罗马人的语感中相邻）。同根词：\textbf{rigor}（严酷；严密——数学的 \textit{rigorous proof} 严密证明；\textit{rigor mortis} 尸僵，法医术语）、\textbf{rigorous}（严格的；缜密的）。意外的表亲：\textbf{frigid}（\textit{frīgus}“寒”$\to$寒冷的；性冷淡的）——rigid 与 frigid 词形之似源自拉丁人把“僵”与“冷”绑在一起的体感：冻僵了，自然 rigid。英语承其感：\textit{a rigid, frigid reception}（冷冰冰的死板接待）两词连用，词源即修辞。
\end{origin}

\vw{rural}{/ˈrʊərəl/}{adj.}

\sense{adj.}{乡村的，农村的}
\col{rural areas / communities / life}
\ex{The clinic serves a wide rural area of some twenty villages.}
\col{rural--urban migration（城乡迁移，社科高频）}
\ex{Rural--urban migration has transformed the demographics of the coast.}

\begin{gram}{rural / urban / suburban 与“郊”的三分}
\textbf{rural}（乡村的）$\leftrightarrow$ \textbf{urban}（城市的）是社科语篇的标准对立轴；\textbf{suburban}（郊区）居中（\textit{suburb} $<$ \textit{sub-} 近 + \textit{urb} 城）。另一对地道反义是 \textbf{the country / the city}（\textit{in the country} 在乡下，country 此义无冠词不可数，与“国家”义判然两分）。美语 \textbf{rural route}（乡邮路线）、\textbf{rural exodus}（乡村人口外流）是新闻常客。形容词辨析配额：\textit{rustic}（乡土风的，常含质朴之褒或土气之贬）比 rural 更带文学滤镜。
\end{gram}

\begin{origin}
拉丁语 \textit{rūrālis}“\textbf{乡野的}”，出自 \textit{rūs, rūris}“\textbf{开阔的田野、乡村}”（原始印欧语 *\textit{rewə-}“\textbf{敞开的空间}”——与 \textbf{room}（房间？其古义正是“空间、余地”，\textit{make room} 留出空间仍存古义）同根！）。城市词 \textit{urbs}（城市$\to$\textbf{urban, suburb, urbane}“城市风度的$\to$温文尔雅的”）与 \textit{rūs} 构成罗马人的城乡二元：\textbf{rural / rustic} 对 \textbf{urban / urbane}。而 \textbf{rustic}（$<$\textit{rusticus}）与 \textbf{rust}（锈？否——\textbf{rust} 铁锈另有词源 \textit{robigo}；但 \textbf{rustic} 与 \textbf{rural} 同根 \textit{rūs}）提醒我们：欧洲语言的“乡下”总与“粗朴”纠缠，与汉语“乡野”的文化负载何其不同又何其相似。
\end{origin}

\end{document}
